
\documentclass{article}
\usepackage{colortbl}
\usepackage{makecell}
\usepackage{multirow}
\usepackage{supertabular}

\begin{document}

\newcounter{utterance}

\centering \large Interaction Transcript for game `taboo', experiment `high\_en', episode 15 with reflexion{-}Qwen3.5{-}4B.
\vspace{24pt}

{ \footnotesize  \setcounter{utterance}{1}
\setlength{\tabcolsep}{0pt}
\begin{supertabular}{c@{$\;$}|p{.15\linewidth}@{}p{.15\linewidth}p{.15\linewidth}p{.15\linewidth}p{.15\linewidth}p{.15\linewidth}}
    \# & $\;$A & \multicolumn{4}{c}{Game Master} & $\;\:$B\\
    \hline

    \theutterance \stepcounter{utterance}  
    & & \multicolumn{4}{p{0.6\linewidth}}{
        \cellcolor[rgb]{0.9,0.9,0.9}{
            \makecell[{{p{\linewidth}}}]{
                \texttt{\tiny{[P1$\langle$GM]}}
                \texttt{You have attempted this same task before and did not succeed. These are the plans you wrote after those attempts.} \\
\texttt{Trial \#1: New plan: I failed because I provided clues that were too broad and allowed the opponent to guess related words (like "PROTECTOR" and "GUARDIAN") instead of the specific target "guard". I should have immediately narrowed the definition to the specific noun "guard" by mentioning a specific object or action, such as "A person who stands at a door to stop people from entering" or "A metal bar on a door that prevents it from opening." Next time, I will focus on the specific physical object or immediate action associated with the word rather than its general function, ensuring the guess is exactly "guard" and not a synonym.} \\
\\ 
\texttt{Trial \#2: New plan: I failed because I provided clues that were too broad and allowed the opponent to guess related words (like "DOORSTOP" and "DOORKEEPER") instead of the specific target "guard". I should have immediately narrowed the definition to the specific noun "guard" by focusing on a unique, non{-}synonymous characteristic, such as "A person who stands at a door to stop people from entering" or "A metal bar on a door that prevents it from opening." Next time, I will focus on the specific physical object or immediate action associated with the word rather than its general function, ensuring the guess is exactly "guard" and not a synonym.} \\
\\ 
\texttt{Trial \#3: New plan: I failed because I provided clues describing the *function* of the word (stopping people, preventing opening) which led the opponent to guess related words like "DOORMAN" or "DOORSTOP" instead of the specific target "guard". I must avoid describing the role or purpose of the word. Instead, I will focus exclusively on the specific physical object that constitutes the word itself, such as "A metal bar on a door" or "A person standing at a door," ensuring the definition points directly to the noun "guard" without triggering synonyms.} \\
\\ 
\\ 
\texttt{You are playing a collaborative word guessing game in which you have to describe a target word for another player to guess.} \\
\\ 
\texttt{Rules:} \\
\texttt{(a) You have to reply in the form: CLUE: <some text>. Guesses from the other player will start with GUESS.} \\
\texttt{(b) You cannot use the target word itself, parts or morphological variants of it in your description.} \\
\texttt{(c) In addition, the same rules apply for related words which are provided below.} \\
\\ 
\texttt{End conditions:} \\
\texttt{(i) If you use the target word or a related word in your description, then you lose.} \\
\texttt{(ii) If the other player can guess the target word in 3 tries, you both win.} \\
\\ 
\texttt{Let us start.} \\
\\ 
\texttt{This is the target word that you need to describe and that the other player needs to guess:} \\
\\ 
\texttt{guard} \\
\\ 
\texttt{Related words are:} \\
\\ 
\texttt{{-} security} \\
\texttt{{-} protect} \\
\texttt{{-} defend} \\
\\ 
\texttt{Important: You are under time pressure, give short descriptions that are to the point!} \\
            }
        }
    }
    & \\ \\

    \theutterance \stepcounter{utterance}  
    & \multicolumn{4}{p{0.6\linewidth}}{
        \cellcolor[rgb]{0.9,0.9,0.9}{
            \makecell[{{p{\linewidth}}}]{
                \texttt{\tiny{[P1$\rangle$GM]}}
                \texttt{CLUE: A metal bar on a door that prevents it from opening.} \\
            }
        }
    }
    & & \\ \\

    \theutterance \stepcounter{utterance}  
    & & & \multicolumn{2}{p{0.3\linewidth}}{
        \cellcolor[rgb]{0.9,0.9,0.9}{
            \makecell[{{p{\linewidth}}}]{
                \texttt{\tiny{[GM$|$GM]}}
                \texttt{continue} \\
            }
        }
    }
    & & \\ \\

    \theutterance \stepcounter{utterance}  
    & & & \multicolumn{2}{p{0.3\linewidth}}{
        \cellcolor[rgb]{0.9,0.9,0.9}{
            \makecell[{{p{\linewidth}}}]{
                \texttt{\tiny{[GM$|$GM]}}
                \texttt{a metal bar on a door that prevents it from opening} \\
            }
        }
    }
    & & \\ \\

    \theutterance \stepcounter{utterance}  
    & & \multicolumn{4}{p{0.6\linewidth}}{
        \cellcolor[rgb]{0.9,0.9,0.9}{
            \makecell[{{p{\linewidth}}}]{
                \texttt{\tiny{[GM$\rangle$P2]}}
                \texttt{You have attempted this same task before and did not succeed. These are the plans you wrote after those attempts.} \\
\texttt{Trial \#1: New plan: I failed to narrow down the specific role despite the hint mentioning "police officer" and "soldier." Next time, I will immediately use the hint to select a specific job title from the suggested examples (e.g., POLICE, SOLDIER, PARENT) rather than guessing broader categories. If the hint suggests multiple roles, I will pick the most common one first to maximize the chance of a correct guess within the three attempts.} \\
\\ 
\texttt{Trial \#2: New plan: I failed to recognize that the target word is a single noun referring to a specific object, instead of a role or a broader category. In the first attempt, I guessed "DOORSTOP" based on the first clue, but the second clue explicitly defined the target as a "person," which I ignored. I then repeated the first guess instead of pivoting to a new word that fits the "person" description. Next time, I will immediately analyze the current clue to determine if it describes a person, object, or action, and select a specific, single{-}word noun that matches that definition. If the clue describes a person, I will guess a specific job title (e.g., DOORMAN, DOORSMITH) rather than a generic role or an object. I will avoid repeating previous guesses and ensure each guess directly addresses the specific constraint of the current clue.} \\
\\ 
\texttt{Trial \#3: New plan: I failed to recognize that the target word is a specific type of person (a "person who stands at a door") rather than a generic role or an object. In the second attempt, I ignored the clue describing a person and guessed an object ("DOORSTOP"). In the third attempt, I guessed a specific job title ("DOORSMITH") but missed the target because the clue was identical to the first one, implying the target is a more common or specific term for that role (likely "DOORKEEPER"). Next time, I will immediately analyze the clue to determine if it describes a person, object, or action. If it describes a person, I will guess a specific, single{-}word noun for that role (e.g., DOORKEEPER, SECURITY, WARDEN) rather than a generic term or an object. I will ensure each guess is a unique, specific word that directly matches the current clue's definition.} \\
\\ 
\\ 
\texttt{You are playing a collaborative word guessing game in which you have to guess a target word that another player describes to you.} \\
\\ 
\texttt{You can make one guess at each trial. You win when you guess the target word. You lose when you cannot guess it in 3 tries.} \\
\\ 
\texttt{After each trial you will get a new hint from the other player which starts with CLUE.} \\
\\ 
\texttt{Make your guesses by just saying the word using the following form: GUESS: <a word>} \\
\\ 
\texttt{Let us start.} \\
\\ 
\\ 
\texttt{CLUE: A metal bar on a door that prevents it from opening.} \\
            }
        }
    }
    & \\ \\

    \theutterance \stepcounter{utterance}  
    & & & \multicolumn{4}{p{0.6\linewidth}}{
        \cellcolor[rgb]{0.9,0.9,0.9}{
            \makecell[{{p{\linewidth}}}]{
                \texttt{\tiny{[GM$\langle$P2]}}
                \texttt{GUESS: DOORSTOP} \\
            }
        }
    }
     \\ \\

    \theutterance \stepcounter{utterance}  
    & & & \multicolumn{2}{p{0.3\linewidth}}{
        \cellcolor[rgb]{0.9,0.9,0.9}{
            \makecell[{{p{\linewidth}}}]{
                \texttt{\tiny{[GM$|$GM]}}
                \texttt{continue} \\
            }
        }
    }
    & & \\ \\

    \theutterance \stepcounter{utterance}  
    & & & \multicolumn{2}{p{0.3\linewidth}}{
        \cellcolor[rgb]{0.9,0.9,0.9}{
            \makecell[{{p{\linewidth}}}]{
                \texttt{\tiny{[GM$|$GM]}}
                \texttt{doorstop} \\
            }
        }
    }
    & & \\ \\

    \theutterance \stepcounter{utterance}  
    & & \multicolumn{4}{p{0.6\linewidth}}{
        \cellcolor[rgb]{0.9,0.9,0.9}{
            \makecell[{{p{\linewidth}}}]{
                \texttt{\tiny{[P1$\langle$GM]}}
                \texttt{GUESS: DOORSTOP} \\
            }
        }
    }
    & \\ \\

    \theutterance \stepcounter{utterance}  
    & \multicolumn{4}{p{0.6\linewidth}}{
        \cellcolor[rgb]{0.9,0.9,0.9}{
            \makecell[{{p{\linewidth}}}]{
                \texttt{\tiny{[P1$\rangle$GM]}}
                \texttt{CLUE: A person who stands at a door to stop people from entering.} \\
            }
        }
    }
    & & \\ \\

    \theutterance \stepcounter{utterance}  
    & & & \multicolumn{2}{p{0.3\linewidth}}{
        \cellcolor[rgb]{0.9,0.9,0.9}{
            \makecell[{{p{\linewidth}}}]{
                \texttt{\tiny{[GM$|$GM]}}
                \texttt{continue} \\
            }
        }
    }
    & & \\ \\

    \theutterance \stepcounter{utterance}  
    & & & \multicolumn{2}{p{0.3\linewidth}}{
        \cellcolor[rgb]{0.9,0.9,0.9}{
            \makecell[{{p{\linewidth}}}]{
                \texttt{\tiny{[GM$|$GM]}}
                \texttt{a person who stands at a door to stop people from entering} \\
            }
        }
    }
    & & \\ \\

    \theutterance \stepcounter{utterance}  
    & & \multicolumn{4}{p{0.6\linewidth}}{
        \cellcolor[rgb]{0.9,0.9,0.9}{
            \makecell[{{p{\linewidth}}}]{
                \texttt{\tiny{[GM$\rangle$P2]}}
                \texttt{CLUE: A person who stands at a door to stop people from entering.} \\
            }
        }
    }
    & \\ \\

    \theutterance \stepcounter{utterance}  
    & & & \multicolumn{4}{p{0.6\linewidth}}{
        \cellcolor[rgb]{0.9,0.9,0.9}{
            \makecell[{{p{\linewidth}}}]{
                \texttt{\tiny{[GM$\langle$P2]}}
                \texttt{GUESS: DOORKEEPER} \\
            }
        }
    }
     \\ \\

    \theutterance \stepcounter{utterance}  
    & & & \multicolumn{2}{p{0.3\linewidth}}{
        \cellcolor[rgb]{0.9,0.9,0.9}{
            \makecell[{{p{\linewidth}}}]{
                \texttt{\tiny{[GM$|$GM]}}
                \texttt{continue} \\
            }
        }
    }
    & & \\ \\

    \theutterance \stepcounter{utterance}  
    & & & \multicolumn{2}{p{0.3\linewidth}}{
        \cellcolor[rgb]{0.9,0.9,0.9}{
            \makecell[{{p{\linewidth}}}]{
                \texttt{\tiny{[GM$|$GM]}}
                \texttt{doorkeeper} \\
            }
        }
    }
    & & \\ \\

    \theutterance \stepcounter{utterance}  
    & & \multicolumn{4}{p{0.6\linewidth}}{
        \cellcolor[rgb]{0.9,0.9,0.9}{
            \makecell[{{p{\linewidth}}}]{
                \texttt{\tiny{[P1$\langle$GM]}}
                \texttt{GUESS: DOORKEEPER} \\
            }
        }
    }
    & \\ \\

    \theutterance \stepcounter{utterance}  
    & \multicolumn{4}{p{0.6\linewidth}}{
        \cellcolor[rgb]{0.9,0.9,0.9}{
            \makecell[{{p{\linewidth}}}]{
                \texttt{\tiny{[P1$\rangle$GM]}}
                \texttt{CLUE: A specific type of guard dog often used for protection.} \\
            }
        }
    }
    & & \\ \\

    \theutterance \stepcounter{utterance}  
    & & & \multicolumn{2}{p{0.3\linewidth}}{
        \cellcolor[rgb]{0.9,0.9,0.9}{
            \makecell[{{p{\linewidth}}}]{
                \texttt{\tiny{[GM$|$GM]}}
                \texttt{continue} \\
            }
        }
    }
    & & \\ \\

    \theutterance \stepcounter{utterance}  
    & & & \multicolumn{2}{p{0.3\linewidth}}{
        \cellcolor[rgb]{0.9,0.9,0.9}{
            \makecell[{{p{\linewidth}}}]{
                \texttt{\tiny{[GM$|$GM]}}
                \texttt{Target word 'guard' (stem=guard) is similar to clue word 'guard' (stem=guard)} \\
            }
        }
    }
    & & \\ \\

    \theutterance \stepcounter{utterance}  
    & \multicolumn{4}{p{0.6\linewidth}}{
        \cellcolor[rgb]{0.95,0.95,0.95}{
            \makecell[{{p{\linewidth}}}]{
                \texttt{\tiny{[P2$\rightarrow$P1]}}
                \texttt{Post{-}episode Reflexion (not sent to GM):} \\
\texttt{New plan: I failed because I described the *function* or *role* of the word (e.g., "prevents it from opening," "stands at a door"), which allowed the opponent to guess related words like "doorstop" or "doorkeeper" instead of the specific target "guard." I must avoid defining the word by what it *does*. Instead, I will focus exclusively on the specific *physical object* that the word names, such as "A metal bar on a door" or "A person standing at a door," ensuring the definition points directly to the noun "guard" without triggering synonyms.} \\
            }
        }
    }
    & & \\ \\

    \theutterance \stepcounter{utterance}  
    & & & \multicolumn{4}{p{0.6\linewidth}}{
        \cellcolor[rgb]{0.95,0.95,0.95}{
            \makecell[{{p{\linewidth}}}]{
                \texttt{\tiny{[P1$\rightarrow$P2]}}
                \texttt{Post{-}episode Reflexion (not sent to GM):} \\
\texttt{New plan: I failed to recognize that the target word is a specific, single{-}word noun for a person who stands at a door, rather than a generic role or a compound word. In the second attempt, I guessed "DOORKEEPER," which is a valid answer, but the game ended in a loss, implying the target was a different, more specific or common single word (likely "DOORMAN" or "WARDEN"). Next time, I will immediately analyze the clue to determine if it describes a person, object, or action. If it describes a person, I will guess a specific, single{-}word noun for that role (e.g., DOORMAN, WARDEN, SECURITY) rather than a generic term, a compound word like "DOORKEEPER," or an object. I will ensure each guess is a unique, specific word that directly matches the current clue's definition.} \\
            }
        }
    }
     \\ \\

\end{supertabular}
}

\end{document}
